\documentclass{jsarticle}
\usepackage{amsmath}
\usepackage{amssymb}
\begin{document}
\title{成りたい自分に近づくために}
\author{}
\date{}
\maketitle

    山口雅旭と申します。原本で既に３冊持っています。彩さんが、「この人と
一緒だとどんなことにもやっていける。」と言っていたのを
本で知ったことについて、私もこの人みたいにはなれないと
思うのと、以前に、彩さんに北川悠仁さん
が、彩さんの父のお墓参りのときに、「お父さんとの約束をやっと果たせた。」
と言っているのを、「聞く、笑う、ツナグ」の本で知ったことを、北川悠仁さん
は、男と男の約束を義徳としてやり遂げたことを、私の今までの言動を鑑みて
私には無理と思うのと、石岡芳隆先生が、性格は直せないが、症状は直せると
言っていたのを、結婚したあとに、夫婦喧嘩をすると子どもにも悪いし、
嫌な思い出夫婦生活をしなければならないので、今頑張れるかどうかです。
彩さんが母親に、「アナウンサーは出来るかどうか」と言われて、やっていける
と、フジテレビに配属されて、私から見ても、木星の霊合みたいな人だなあと、
私も中学校のときの河合塾と２０１５年の剣道とベッセルの就職も、
周りに対する気持ちと、この負けず嫌いもいいとも思い、私も
この体験をしていたなと思いました。


    速読法には音読以上の記憶力を授かる
速読があるのと、これとは別の能力の睡眠力を、SRSのときと中国新聞の記事を毎日読んでいた幼少期にも、
この睡眠力を速読と音読法とで最近知りました。私も初めはすごく睡魔に
バタンキューとする睡眠力が速読法を
したときの後にありましたが、睡魔を我慢していたら、この能力が消えて、
短眠だけになりました。ちゃんと長時間寝ないと病気は治りにくいのも、速読法
での、睡眠力と別の記憶力の能力も体験していました。　
寝るのを我慢してはいけないと先生達から学びました。
いつまでも覚えていてはだめらしいのがわかり、９時間寝れば覚えていなくていい
状態になるらしいのも教えられました。結局長時間決まった時間に寝ないといけないのも教えられました。アナウンサーがどれだけ睡眠法を駆使しているかも、
めざましテレビの皆藤愛子さんと中野美奈子さんと彩さんで以前から聞いていました。
北野武さんの話も聞いていて、決まった時間が大切なのも教えられました。看護師さんの夜勤も北野武さんと同じく、決まった時間に睡眠をとるのも知りました。　
社会的な規則があるのも知りました。
みんな速読法をしている上に、睡眠を大切にしているのも、SRS速読法のレム睡眠
とノンレム睡眠の浅瀬の海で快眠するのと、予知がある深い睡眠の深海に
潜って熟睡する睡眠法があるのも最近になって知りました。
 

   母親が順子お姉さんに料理の本を
渡していたのを、この本を私は
勉強して、ベッセルの退職したあとに、SRSのもぐらたたきの練習で得た触覚能力と
川島隆太先生からの、記憶力で目分量と具材が煮えているかを触覚からわかるように
なり、やれば出来る体験をしていました。小林カツ代さんの
本で、最近になって、この料理家はマクロビの欠点を知っているらしく、
お肉と魚をうまい具合に野菜と組み合わせているのに気づきました。
塩分濃度を気をつけていたら大丈夫らしいのも、この人から学びました。
初めのときは目分量を時間と直感からやり始めて、この感覚
で、後になって、目分量の目安がわかり、これと同時に具材もわかるようになり、
レシピを見なくても、自分のオリジナル料理
としての混ぜ料理が出来るようになりました。
２００６年の１０月に彩さんが　子どもは３人欲しいと、
悠仁さんに言っているのを、私も兄弟３人いて、幼少期のときに兄弟喧嘩していても、
私が高等学校のときに、姉から数学とZ会を教えてもらえたことで、富山大学に
合格できて、その上に、姉からピアノも教えてもらえて、NHKのおかげもあり
、作曲までも出来るようになりました。
平成１８年の２００６年河合塾の３月から６月まで通っていたときに、
圭一さんが、YMCAのときも、このときの圭一さんがいた。２００３年
のときも圭一さんがいた。YMCAのときには、このときの圭一さんが、
私がああ、女神様を読んでいたのを最近になって知っていたのに
気づいた。薬学とこの主人公の蛍一さんが、工学系の専門家であり、
バイクが好きで、すごくラッキーな性格であり、薬学はウルドで、
すごくこの蛍一さんとベルダンディーさんが好きで読んでいたのを
覚えている。傾倒して読んでいた。富山大学のときは、
このときの圭一さんが運輸会社に勤めたいと言って、自動車学校に通っていた。
ハローワークが最後に、まとめていた。
２００６年の河合塾は、彩さん関係の人に助けられていた。
大塚さんが、彩さんが休暇を取っていたときに、こたつで見ていたでしょうと
言っていたのを、彩さんがありがとうございます。と言っていたのを
最近になって気づいた。全部は、彩さんの「聞く、笑う、ツナグ」
を読んでわかった。
彩さんも２０１１年に北川悠仁さんと御祓神社で挙式を挙げているのを
読んで、自分もこの２００６年の９月に御祓のために、杉の桶で朝に御祓をしていて
縁側も開けようとしていたのを、この御祓は冷たいし寒かったのを思い起こしました。
このときに、夕方に散歩で、ホンダのお店の近くで、遠くから、フルメタの
テレサ・テスタロッサさんの制服を着た人が来ているのを、最近になって気づきました。
２００５年の６月に初めてSRS速読法が、手の開閉運動と屈伸運動で音が文字から
外れて、速読法が出来るようになりました。指回し体操も速読の大の目と周の目、
中の目と確の目と出来るようになり、新聞社の息子さんの影響でもあり、作業療法士さん
の言葉でもあり、この年はいろんな体験をしました。
２０１０年のベッセルのニューキャッスルホテルでのビアホールのアルバイトで、
坂本孝文シェフのアドバイスの、「枚数を数えずに数えること」
が、最近の速読技法での、一定でのブロックを単位とする、「５ブロックごとに
単位として」一つずつは数えずに、あるまとまりでどんどんカウントすること、
１２ブロックをひとかたまりに数えることで、この技法が出来ている。この技法は
音を使わずに数えることがポイントでもあるのは、
坂本シェフに言われたときに、速読だとはわかっていたが、やれる状態ではなかった。
岡崎茜さんは直ぐにステーキを焼くのも４切れを一片に掴んで鉄板に並べて
この技法を既にやっていた。私は最近になって、いろんなものを数えるのにこの技を
練習しています。
「聞く、笑う、ツナグ」を読んで彩さんが、「読んだ内容を自分の言葉として
書き出すといいですよ。」という言葉のアドバイスを受けて、今初めての論文を
書いている最中でもあります。実際書き出しが今までの数学の分野の領域が
初めて書き出しとして出来つつある次第でもあります。
アカシックレコードを知ってから、この場所にアクセスしたら、
本当に数学のアカシックレコードがあり、実際に読んでもいます。
書き出しも出来てもいます。
彩さんと同じかどうかはわかりませんが、私も昔は負けず嫌いでもありました。
２０１５年の剣道は、トキ子おばあちゃんに助けられていたのを、私も他力の
力を大切にと、剣道をしてよかったかを後悔とまではですが、いろんな経験を
しています。彩さんにまた会えたらいいです。　　　
フジテレビのるろうに剣心で昔は剣道が嫌だったのが好きになり、
剣友会も剣心の垂れ幕を作っていました。
ど根性精神は、苫米地博士は推奨していないですが、目的によって使い分ける
ことも功利としては手ともおもうかもしれないです。スコトーマにもなっていると
思う次第でもあります。
栗田博士も私に２０１０年の１０月に電話で「またどこかで会えたら素敵ですね。」
と言っていました。
彩さんは苫米地博士の研究を知っているらしく、未練を穫るとも言っているし、
北川悠仁さんと、本当かどうか疑わしい限りでもありますが、私が北川悠仁さん
にはなれないのに付き合ってくれているらしく、これが本当かどうか、
本当だったら嬉しいの一言でもあり、１０年以上前にこのことを知っていたら、
どんだけよかったか、一言で言い尽くせない思いでもあります。
時間は未来から過去も私も博士を知ってから、物理の分野でゼータ関数としてこの
知識を論文に取り入れていたら、全分野がこの時間の仕組みを内在しているのを、
理路整然と論文に書き出している最中でもあります。
すべては、彩さんの本からの知識の吸収でまかり通っている次第でもあります。
フジテレビの伊藤清社長に私は叱られているのを、すごく反省しています。
大塚範一さんにも言われているのもわかっています。
彩さんに謝るのは、このことです。
めざましテレビにも言われているのもわかっています。
「聞く、笑う、ツナグ」を読んで、すごくわかっています。
今、治療の経過でもあります。
瞑想法は大変でもあります。
２００６年の９月４日は、大塚範一さんに軽部真一さん、吉岡さん、内田理沙さん、
皆藤愛子さん、中野美奈子さん、伊藤利尋さん、林さん、めざましテレビの
吉岡さんは、好きです、このときの彩さんが一緒に
「屋久どんさん、めざましテレビみてください。
こちそうございました。」と鹿児島に行って、サインの色紙を
インターネットで見ていました。めざましテレビの
メンバーを代表に選んで、車で母親と福山市に来ているのを
後から知りました。何日かあとに彩さんが車のドアーに顔を
隠しているのを２回見ています。母親には会っていて、
なぜかお兄さんが２回見ている気がします。一回目は
バスケットボールでのことで、２回目は、福祉専門学校の
近くでもあります。統合失調症は、本当に大変でもあり、
芸能人と付き合っている妄想が、これは、言葉上と実体験からの症状でもあり、
大変なのであり、私のことは、勘違いもあります。　
双極性のうつは、占いを知って、緑生の特質の人がかかる病気であるのを、
広大の清水さんと高橋さん、中村さん、中野さん、惣引さん、
芸術では、このうつにメルエールの病を合わせると、色彩がなんとも言えない、
その上に風景が躍動感の渦巻を描いている月の絵と印象的な　
魅力の絵を描くゴッホのひまわりの絵画と、アインシュタイン博士の
分裂病的気質と、病気が良い影響を与えるのも緑生の能力と、３２、３３野
の優しさと言えるのも、１９９１年からの美鈴ケ丘の塾からの学問の勉強から
はじまり、現在の知識からわかったことでもあります。病気から学んだことでも
あります。相当に、２０１０年は大変な人間関係であったです。　
福山市のビオーネが名産と知っていたら、ちゃんと皆藤愛子さんの話を
聞いていたら、とっくに気づけました。本当に２０１０年と２０１１年は、
迷惑をかけました。微分幾何の量子化が本当だったら、安心して産んでもいいと
思います。彩さんが、２０１２年に一回流産していた後に、罪悪感から布団に潜り込み、
北川悠仁さんも大変だったらしいと、「彩日記」で読んでいます。
最近になって、「彩日記」から結婚線が一つらしいと気づきました。
決心して出産に立ち寄ったことを、本当に好きなのかーと思いました。
私は同じ長さで２本線だから、同じ人だなとも思いました。
私の結婚線はしつこいくらい長いです。好きな人がいるのをちゃんと大切に
しないと、本当は一本線で終わっていたのかとも思っています。
２０１１年は本当にしかたなかったです。
彩さんの、「育児日記」を読んで、彩子さんが「赤ちゃんにだけなりたい」と
関係副詞の係助詞で「にだけ」と２歳で言っているのを、
彩子さんが言っていて、
美穂さんが「お父さんの真似」としているインスタグラムから、
最近は、私も「アカシックレコード」から、微分幾何の量子化でガンマ関数とゼータ関数をJones多項式からオイラーの定数と数学の数値が生まれてきた数式が大域的微分多様体
となって、恵みをもらいました。全部は彩さんの本のラブレターから得た恩恵でした。
私は本当に本が好きなのです。なぜか、親に洗脳されている現在でもあります。

\end{document}

